\chapter{Introduction}
\label{ch:intro}

This chapter presents an overview of the proposed research on coordinated
propaganda detection in Bangla digital media. It outlines the project background,
problem statement, core motivation, research objectives, methodology, expected
outcomes, and the structural organization of the report.

\section{Project Overview}
Bangladesh's digital media ecosystem has become a major source of information for
millions of people, alongside continued growth in internet access and social media
use. According to DataReportal's \emph{Digital 2025: Bangladesh}, the country had
\textbf{77.7 million internet users} and \textbf{60.0 million social media user
identities} in early 2025, demonstrating how strongly online platforms now shape
news circulation and public opinion.

Within this environment, \textbf{propaganda} refers to strategically crafted
information designed to influence public attitudes, beliefs, or actions. In the
digital age, such content is often spread through coordinated campaigns, repeated
messaging, and networked amplification, making it more scalable and harder to
detect than isolated misinformation. This issue is especially important in
Bangladesh, where online platforms are deeply embedded in political communication,
election-related discourse, and reactions to controversial events, all of which
increase the potential impact of coordinated influence operations~\cite{hristakieva2022}.

\subsection*{Problem Statement}
Most automated propaganda detection systems to date have been developed for English
and other languages with abundant digital resources. Bangla, by contrast, still
lacks large annotated datasets and detection frameworks capable of identifying
coordinated propaganda across text, images, and cross-platform activity, which
leaves a substantial research and practical gap~\cite{hossain2020}. This project
responds to that gap by proposing a Bangla propaganda dataset and an AI-based
detection framework tailored to Bangladesh's online media environment. The goal is
to support more accurate, scalable, and context-sensitive detection of coordinated
propaganda campaigns in Bangla-language digital spaces~\cite{hossain2020,kamruzzaman2023}.

\section{Motivation}
The motivation for this project extends beyond building a detection system. In an
increasingly digital society, people rely on online information to understand
current events, form opinions, and make decisions that affect their communities and
democratic processes. When coordinated propaganda manipulates information at scale,
it can undermine public trust, distort discourse, and reduce people's ability to
make informed judgments~\cite{hristakieva2022}.

This project is grounded in the belief that individuals should be able to think
independently rather than have their perceptions shaped by coordinated
manipulation. Supporting that independence requires a digital environment in which
misleading influence campaigns can be identified and examined more effectively.
Technology alone cannot eliminate propaganda, but it can provide tools that improve
transparency and help users evaluate online content more critically~\cite{ng2023}.

From a research standpoint, the scarcity of Bangla-language resources for this task
presents an opportunity. Existing Bangla work has already shown the value of
annotated datasets and benchmark systems for low-resource language processing, and
the same need extends to multimodal and propaganda-oriented settings. In this way,
the project aims to contribute both to advances in AI research and to a more
trustworthy online information ecosystem in Bangladesh~\cite{hossain2020}.

\section{Objectives}
The primary purpose of this research is to build a foundational AI framework and
dataset for the Bangla digital media environment that can effectively detect and
analyze coordinated propaganda campaigns~\cite{ng2023}.

As this research is currently in its initial phase, the specific objectives are
enumerated below:
\begin{enumerate}
  \item \textbf{Create a large-scale, multimodal Bangla dataset:} Compile and
        annotate a comprehensive dataset spanning diverse domains, such as politics
        and news, including both textual and visual elements, using hybrid
        annotation pipelines that combine human expertise with LLM assistance to
        scale the process efficiently~\cite{hossain2020}.
  \item \textbf{Develop an AI model for propaganda detection:} Architect and train
        a multimodal machine learning framework optimized for the linguistic
        nuances of Bangla to classify manipulative and propagandistic
        content~\cite{hossain2020}.
  \item \textbf{Analyze coordinated campaigns:} Combine semantic, temporal, and
        network-level signals to detect organized, scalable influence operations
        rather than isolated pieces of misinformation~\cite{ng2023}.
  \item \textbf{Evaluate model performance:} Conduct baseline evaluations against
        current state-of-the-art models, such as fine-tuned transformers and
        cross-lingual models, to assess accuracy, scalability, and
        context-sensitivity in the Bangladeshi digital landscape~\cite{hossain2020}.
  \item \textbf{Establish open resources for future research:} Make the annotated
        dataset and architectural findings accessible to the academic community to
        support further advances in NLP, multimodal machine learning, and digital
        media analysis for low-resource languages~\cite{hossain2020}.
\end{enumerate}

\section{Methodology}
\label{sec:method-overview}
To achieve these objectives, the research follows a structured, multi-phase
methodology. Figure~\ref{fig:workflow} presents the complete pipeline at a glance:
raw Bangla content is ingested and normalized, passed through a hybrid
LLM-plus-human annotation stage to yield a verified multimodal dataset, and then
processed by three parallel analysis streams (text, visual, and network) whose
outputs are combined by a cross-modal campaign-fusion head and delivered through a
web dashboard and monitoring API.
\begin{figure}[H]
\centering
\begin{adjustbox}{max width=\textwidth,max height=0.85\textheight}
\begin{tikzpicture}[
  wfn/.style={rectangle,draw=black!70,rounded corners=2pt,align=center,
              font=\small,minimum height=13mm,text width=34mm,inner sep=4pt},
  wfsrc/.style ={wfn,fill=gray!10},
  wfproc/.style={wfn,fill=blue!7,minimum height=20mm},
  wfsto/.style ={wfn,fill=green!8},
  wfout/.style ={wfn,fill=orange!10,text width=55mm},
  stl/.style={rotate=90,anchor=south,font=\scriptsize\bfseries,align=center,text width=30mm},
  wfgrp/.style={draw=black!45,dashed,rounded corners=4pt,inner sep=8pt}
]
% ---- Stage 1
\node[wfsrc] (src)   at (-4.6,0)    {\textbf{Bangla Digital Media Sources}\\[2pt]{\scriptsize Facebook \quad Twitter/X \quad News Outlets}};
\node[wfsrc] (clean) at (0,0)       {\textbf{Data Cleaning}\\[2pt]{\scriptsize Unicode Normalization}};
\node[wfsrc] (annot) at (4.6,0)     {\textbf{Hybrid Annotation Pipeline}\\[2pt]{\scriptsize LLM Pre-Annotation\\+ Human Verification}};
\node[wfsto] (ds)    at (0,-2.4)    {\textbf{Verified Multimodal Dataset}};
\draw[flow] (src) -- (clean);
\draw[flow] (clean) -- (annot);
\draw[flow] (annot.south) |- (ds.east);
\node[wfgrp,fit=(src)(clean)(annot)(ds)] (st1) {};
\node[stl] at (st1.west) {1.\ Ingestion \& Preprocessing};

% ---- Stage 2
\node[wfproc] (txt) at (-4.6,-6.2) {\textbf{Text Stream}\\[2pt]{\scriptsize BanglaBERT / XLM-R\\Semantic \& Persuasive Cue Modeling}};
\node[wfproc] (vis) at (0,-6.2)    {\textbf{Visual Stream}\\[2pt]{\scriptsize ResNet / ViT\\Image--Text Alignment \& Meme Analysis}};
\node[wfproc] (net) at (4.6,-6.2)  {\textbf{Network Stream}\\[2pt]{\scriptsize Temporal Clustering\\Coordinated Account Interaction Graphs}};
\node[wfgrp,fit=(txt)(vis)(net)] (st2) {};
\node[stl] at (st2.west) {2.\ Multimodal AI Framework};

% bus: dataset -> three streams
\draw[line] (ds.south) -- (0,-4.4);
\draw[line] (-4.6,-4.4) -- (4.6,-4.4);
\draw[flow] (-4.6,-4.4) -- (txt.north);
\draw[flow] (0,-4.4)    -- (vis.north);
\draw[flow] (4.6,-4.4)  -- (net.north);

% ---- Stage 3
\node[wfout] (fuse) at (0,-9.8)  {\textbf{Cross-Modal Campaign\\Fusion Head}};
\node[wfout] (out)  at (0,-11.7) {\textbf{Propaganda Category}\\{\scriptsize Confidence \& Technique}};
\node[wfout] (dash) at (0,-13.6) {\textbf{Web Dashboard}\\{\scriptsize \& Monitoring API}};
\draw[flow] (fuse) -- (out);
\draw[flow] (out) -- (dash);
\node[wfgrp,fit=(fuse)(out)(dash)] (st3) {};
\node[stl] at (st3.west) {3.\ Synthesis \& User Deployment};

% bus: three streams -> fusion
\draw[line] (txt.south) -- (-4.6,-8.4) -- (0,-8.4);
\draw[line] (net.south) -- (4.6,-8.4)  -- (0,-8.4);
\draw[flow] (vis.south) -- (fuse.north);
\end{tikzpicture}
\end{adjustbox}
\caption{Overall workflow of the proposed Bangla propaganda detection framework.}
\label{fig:workflow}
\end{figure}

The workflow proceeds through five core phases~\cite{ng2023}:
\begin{itemize}
  \item \textbf{Data Collection and Integration:} Raw data containing text and
        visual elements across politics, sports, and entertainment is
        systematically gathered using automated web scrapers and validated
        manually.
  \item \textbf{Data Cleaning and Preprocessing:} Unstructured text is cleaned,
        standardized, tokenized, and normalized into machine-readable formats
        tailored for Bangla NLP tasks.
  \item \textbf{Dataset Annotation and Construction:} A hybrid pipeline uses LLMs
        for initial pre-annotation and span extraction, followed by human expert
        verification to ensure label quality and reduce
        bias~\cite{hossain2020}.
  \item \textbf{Multimodal Model Development:} Semantic features, temporal metrics,
        and social network structure are fused to pinpoint coordinated campaigns
        rather than standalone posts~\cite{ng2023}.
  \item \textbf{Evaluation and Comparative Analysis:} Benchmarks are run against
        transformer architectures and cross-lingual models to measure real-world
        performance on Bangladeshi digital media contexts~\cite{hossain2020}.
\end{itemize}

\section{Project Outcome}
By the end of this phase of the research, the project is expected to deliver the
following major concrete outputs~\cite{hossain2020}:
\begin{itemize}
  \item An annotated multimodal Bangla propaganda dataset covering multiple public
        domains.
  \item A trained AI framework designed for detecting coordinated digital
        propaganda.
  \item A comprehensive baseline evaluation comparing the proposed architecture
        against state-of-the-art models.
  \item Openly available dataset resources and benchmarks to foster future NLP and
        multimodal research for under-resourced languages~\cite{hristakieva2022}.
\end{itemize}

\section{Organization of the Report}
The remainder of this report is organized as follows:
\begin{itemize}
  \item \textbf{Chapter~\ref{ch:background} (Background):} Presents the necessary
        preliminaries, reviews similar applications and related research, and
        identifies the research gaps relevant to coordinated propaganda detection
        in Bangla online media.
  \item \textbf{Chapter~\ref{ch:design} (Project Design):} Describes the project
        requirements, context diagram, data flow diagrams, user interface design,
        detailed methodology, project plan, and task allocation.
  \item \textbf{Chapter~\ref{ch:standards} (Standards and Design Constraints):}
        Discusses the standards followed, realistic design constraints, cost
        analysis, and the mapping of the project to complex engineering problem
        solving and activity criteria.
\end{itemize}